\documentclass[UTF8,a4paper,12pt,oneside,openany,fontset=fandol]{ctexbook}
\usepackage{amsmath,amssymb,graphicx,geometry,fancyhdr,enumitem,array}
\usepackage{empheq}
\usepackage{tikz}
\usetikzlibrary{arrows.meta,calc,patterns,decorations.markings}
\tikzset{ray/.style={thin,postaction={decorate},decoration={markings,mark=at position .6 with {\arrow{Stealth}}}}}
\usepackage[unicode,hidelinks]{hyperref}
\geometry{left=23mm,right=23mm,top=22mm,bottom=22mm,headheight=16pt}
\pagestyle{fancy}\fancyhf{}
\fancyhead[L]{\small 赵凯华新概念光学做题本}
\fancyhead[R]{\small\nouppercase{\leftmark}}\fancyfoot[C]{\thepage}
\renewcommand{\chaptermark}[1]{\markboth{第\thechapter 章\quad #1}{}}
\setlength{\parindent}{2em}\setlength{\emergencystretch}{3em}
\linespread{1.12}
\ctexset{chapter={beforeskip=10pt,afterskip=20pt}}
\setlist[enumerate]{label=（\arabic*）,leftmargin=2.5em,itemsep=2pt,topsep=3pt}
\newcommand{\problem}[2]{\par\phantomsection\pdfbookmark[1]{习题 #1}{problem-#1}\noindent\textbf{习题 #1}\quad #2\par}
\newsavebox{\figbox}
% Compact diagrams reserve space for handwritten work. Dense multi-part
% diagrams have a separate height cap so their labels remain readable.
\newlength{\maxfigheight}
\foreach\figid in {1_18,2_10,3_16,4_5,book_3_33,book_7_9}{\expandafter\gdef\csname detailedfig@\figid\endcsname{1}}
\newcommand{\drawfig}[2]{\begingroup\setlength{\maxfigheight}{38mm}\ifcsname detailedfig@#2\endcsname\setlength{\maxfigheight}{48mm}\fi\sbox{\figbox}{\resizebox{#1\linewidth}{!}{\input{tikzfigs/fig_#2.tex}}}\sbox{\figbox}{\scalebox{0.9}{\usebox{\figbox}}}\ifdim\ht\figbox>\maxfigheight\resizebox{!}{\maxfigheight}{\usebox{\figbox}}\else\usebox{\figbox}\fi\endgroup}
\def\figlabel#1_#2\relax{#1-#2}
\newcommand{\fig}[2][0.46]{\par\addvspace{4pt}{\centering\drawfig{#1}{#2}\par\vspace{2pt}{\footnotesize 习题\figlabel#2\relax\ 图}\par}\addvspace{3pt}}
\newcommand{\bookfig}[3][0.8]{\par\addvspace{4pt}{\centering\drawfig{#1}{#2}\par\vspace{2pt}{\footnotesize #3}\par}\addvspace{3pt}}
% A short question and its diagram form one aligned, non-floating block.
\newcommand{\problemfig}[4][0.35]{\par\phantomsection\pdfbookmark[1]{习题 #2}{problem-#2}\noindent\begin{minipage}[t]{\dimexpr\linewidth-#1\linewidth-6mm\relax}\vspace{0pt}\textbf{习题 #2}\quad #3\end{minipage}\hfill\begin{minipage}[t]{#1\linewidth}\vspace{0pt}\centering\drawfig{1}{#4}\par\vspace{3pt}{\footnotesize 习题\figlabel#4\relax\ 图}\end{minipage}\par}
\newcommand{\lenstable}{\begin{center}\small\renewcommand{\arraystretch}{1.6}\begin{tabular}{|c|*{7}{p{1cm}|}}\hline 物距 $s/\mathrm{cm}$&$-24$&$-12$&$-6.0$&0&12&24&36\\\hline 像距 $s'/\mathrm{cm}$&&&&&&&\\\hline 横向放大率 $V$&&&&&&&\\\hline 像的虚实&&&&&&&\\\hline 像的正倒&&&&&&&\\\hline\end{tabular}\end{center}}
\newcommand{\missing}[1]{\par{\small\itshape〔待补资料：#1〕}\par}
\newcommand{\editorial}[1]{\par\smallskip{\small〔校注：#1〕}\par}
\newcommand{\answerarea}{\par\vfill}
\begin{document}
\begin{titlepage}\centering\vspace*{.24\textheight}
{\Huge\bfseries 赵凯华新概念光学\\[1em]做题本\par}
\vspace{2cm}{\Large 习题整理与作答空间\par}\vfill
{\large 姓名：\underline{\hspace{4cm}}\quad 日期：\underline{\hspace{3cm}}\par}
\end{titlepage}
\tableofcontents\clearpage
\chapter{光和光的传播}
\problem{1-1}{太阳与月球的直径分别是 $1.39\times10^6\,\mathrm{km}$ 和 $3.5\times10^3\,\mathrm{km}$，太阳到地面的距离为 $1.50\times10^8\,\mathrm{km}$，月球到地面的距离为 $3.8\times10^5\,\mathrm{km}$。试计算，地面上能见到日全蚀区域的面积（可把该区域的地面视为平面）。}
\answerarea
\problemfig[0.34]{1-2}{由立方体的玻璃切下一角做成的棱镜，称为隅角棱镜（见题图）。证明从斜面射入的光线经其它三面全反射后，出射线的方向总与入射线相反。设想一下，这样的棱镜可以在什么场合发挥作用。}{1_2}
\answerarea
\clearpage
\problemfig[0.38]{1-3}{光线射入如题图所示的棱镜，经两次折射和反射后射出。
\begin{enumerate}\item 证明偏向角与入射方向无关，恒等于 $2\alpha$。\item 在此情况下，能否产生色散？\end{enumerate}}{1_3}
\answerarea
\problem{1-4}{试证明：当一束光线通过平行平面玻璃板时，出射光线方向不变，但产生侧向平移。当入射角 $\theta$ 很小时，位移为
\[x=\frac{n-1}{n}\theta t,\]
式中 $n$ 为玻璃板的折射率，$t$ 为其厚度。}
\answerarea
\clearpage
\problem{1-5}{证明：光线相继经过几个平行分界面的多层介质时，出射光线的方向只与两边的折射率有关，与中间各层介质无关。}
\answerarea
\problem{1-6}{顶角 $\alpha$ 很小的棱镜称为光楔（wedge）。证明光楔使垂直入射的光线产生偏向角 $\delta=(n-1)\alpha$，其中 $n$ 是光楔的折射率。}
\answerarea
\clearpage
\problem{1-7}{本题图所示是一种求折射线方向的简捷作图法。例如为了求光线通过棱镜的路径（图 b），可如图 a 以 $O$ 为中心作二圆弧，半径正比于折射率 $n$、$n'$（设 $n>n'$）。作 $OR$ 平行于入射线 $DE$，作 $RP$ 平行于棱镜第一界面的法线 $N_1N'_1$，则 $OP$ 的方向即为第一次折射后光线 $EF$ 的方向。再作 $QP$ 平行于第二界面的法线 $N_2N'_2$，则 $OQ$ 的方向即为出射线 $FG$ 的方向，从而 $\angle ROQ=\delta$ 为偏向角。试论证此法的依据。\fig[0.82]{1_7}}
\answerarea
\problem{1-8}{利用上题的图证明最小偏向角的存在及（1.23）式。
\begin{equation*}
n=\frac{\sin\dfrac{\alpha+\delta_{\min}}{2}}{\sin\dfrac{\alpha}{2}}.
\tag{1.23}
\end{equation*}
其中 $\alpha$ 为棱镜顶角，$\delta_{\min}$ 为最小偏向角。}
\answerarea
\clearpage
\problem{1-9}{已知棱镜顶角为 $60^\circ$，测得最小偏向角为 $53^\circ14'$，求棱镜的折射率。}
\answerarea
\problem{1-10}{顶角为 $50^\circ$ 的三棱镜的最小偏向角是 $35^\circ$，如果把它浸入水中，最小偏向角等于多少？（水的折射率为 $1.33$。）}
\answerarea
\clearpage
\problemfig[0.32]{1-11}{如本题图所示，在水中有两条平行光线1和2，光线2射到水和平行平板玻璃的分界面上。
\begin{enumerate}\item 两光线射到空气中是否还平行？\item 如果光线1发生全反射，光线2能否进入空气？\end{enumerate}}{1_11}
\answerarea
\problem{1-12}{计算光在下列介质之间穿行时的全反射临界角：（1）从玻璃到空气，（2）从水到空气，（3）从玻璃到水。}
\answerarea
\clearpage
\problem{1-13}{设光导纤维玻璃芯和外套的折射率分别为 $n_1$ 和 $n_2$（$n_1>n_2$），垂直端面外介质的折射率为 $n_0$（见题图）。试证明，能使光线在纤维内发生全反射的入射光束的最大孔径角 $\theta_1$ 满足下式：
\[n_0\sin\theta_1=\sqrt{n_1^2-n_2^2}.\]
（$\theta_1$ 称为纤维的数值孔径。）\fig[0.7]{1_13}}
\answerarea
\problem{1-14}{光导纤维外套由折射率为 $1.52$ 的冕牌玻璃做成，芯线由折射率为 $1.66$ 的火石玻璃做成，求垂直端面的数值孔径。}
\answerarea
\clearpage
\problem{1-15}{极限法测液体折射率的装置如本题图所示，$ABC$ 是直角棱镜，其折射率 $n_g$ 为已知。将待测液体涂一薄层于其上表面 $AB$，覆盖一块毛玻璃，用扩展光源在掠入射的方向照明。从棱镜的 $AC$ 面出射的光线的折射角将有一下限 $i'$（用望远镜观察，则在视场中出现有明显分界线的全明半暗区）。试证明，待测液体的折射率 $n$ 可按下式算出：
\[n=\sqrt{n_g^2-\sin^2 i'}.\]
用这种方法测液体折射率，测量范围受到什么限制？\fig{1_15}}
\answerarea
\problem{1-16}{在空气中钠黄光的波长为 $589.3\,\mathrm{nm}$，问：（1）其频率有多大？（2）在折射率为 $1.52$ 的玻璃中其波长为多少？}
\answerarea
\clearpage
\problem{1-17}{在熔融石英中波长为 $550.0\,\mathrm{nm}$ 的光频率为多少？已知折射率为 $1.460$。}
\answerarea
\problem{1-18}{填充下表中的空白：
\begin{center}\renewcommand{\arraystretch}{1.5}\begin{tabular}{|c|c|c|c|c|}\hline
谱线&\multicolumn{2}{c|}{F线}&\multicolumn{2}{c|}{D线}\\\hline
介质&真空&水&真空&水\\\hline
折射率&&1.337&&1.333\\\hline
波长&$486.1\,\mathrm{nm}$&&$589.3\,\mathrm{nm}$&\\\hline
频率&&&&\\\hline 光速&&&&\\\hline
\end{tabular}\end{center}}
\answerarea
\clearpage
\problem{1-19}{一只船以速度 $v$ 在水面上行驶，$v$ 比水面波的速率 $u$ 大。用惠更斯作图法证明，在水中出现一圆锥形波前，其半顶角 $\theta$ 由下式给出：
\[\sin\theta=\frac{u}{v}.\]
此波称为船波，超音速飞机在空气中产生的冲击波，也是这样产生的。设想一下，若电子以大于介质中光速的速度在介质中作匀速运动，会产生什么现象？}
\answerarea
\problem{1-20}{证明图1-18中光线 $A_1C_1$、$A_2B_2C_2$、$A_3B_3C_3$、$\cdots$、$A_nB_n$ 的光程相等。
\bookfig[0.82]{1_18}{图1-18\quad 用惠更斯原理解释反射定律和折射定律}}
\answerarea
\clearpage
\problem{1-21}{证明图1-18中光线 $A_1D_1$、$A_2B_2D_2$、$A_3B_3D_3$、$\cdots$、$A_nB_n$ 的光程相等。
\bookfig[0.82]{1_18}{图1-18\quad 用惠更斯原理解释反射定律和折射定律}}
\answerarea
\problemfig[0.37]{1-22}{在离桌面 $1.0\,\mathrm m$ 处有盏 $100\,\mathrm{cd}$ 的电灯 $L$，设 $L$ 可看作是各向同性的，求桌上 $A$、$B$ 两点的照度（见本题图）。}{1_22}
\answerarea
\clearpage
\problem{1-23}{若上题中电灯 $L$ 可垂直上下移动，问怎样的高度使 $B$ 点的照度最大。}
\answerarea
\problem{1-24}{（1）设天空为亮度均匀的朗伯体，其亮度为 $B$，试证明，在露天水平面上的照度为 $E=\pi B$；（2）在上面的计算中，与我们假设天空是怎样形状的发光面有无关系？与被照射面的位置有无关系？}
\answerarea
\clearpage
\problem{1-25}{试证明，一个理想漫射体受到照度为 $E$ 的辐射时，反射光的亮度 $B=E/\pi$。}
\answerarea
\problem{1-26}{阳光垂直照射地面时，照度为 $1.0\times10^5\,\mathrm{lx}$。若认为太阳的亮度与光流方向无关，并忽略大气对光的吸收，且已知地球轨道半径为 $1.5\times10^8\,\mathrm{km}$，太阳的直径为 $1.4\times10^6\,\mathrm{km}$，求太阳的亮度。}
\answerarea
\clearpage

\chapter{几何光学成像}
\problem{2-1}{以第一章第1节例题2中所用的光线追迹作图法证明图2-6中 $Q$ 和 $Q'$ 是一对共轭点。
\par\medskip\noindent
\begin{minipage}[t]{.35\linewidth}\vspace{0pt}\centering
\drawfig{1}{book_2_6}\par\smallskip
{\footnotesize 图2-6\quad 折射球面的齐明点}
\end{minipage}\hfill
\begin{minipage}[t]{.61\linewidth}\vspace{0pt}\centering
\drawfig{1}{book_2_19}\par\smallskip
{\footnotesize 图2-19\quad 用作图法求二次成像（补充参考图）}
\end{minipage}\par
\missing{题干引用的第一章第1节例题2正文；图2-6及补充图2-19已补入。}}
\answerarea
\clearpage
\problem{2-2}{证明（2.4）式。图2-6中，角度 $u=\angle MQA$、$u'=\angle MQ'A$ 与距离 $\overline{QA}$、$\overline{AQ'}$ 之间存在下列关系：
\begin{equation*}
\frac{\sin u}{\sin u'}=\frac{\overline{AQ'}}{\overline{QA}}.
\tag{2.4}
\end{equation*}
\bookfig[0.38]{book_2_6}{图2-6\quad 折射球面的齐明点}}
\answerarea
\clearpage
\problem{2-3}{根据反射定律推导球面反射镜的物像距公式（2.23）和焦距公式（2.24）。}
\answerarea
\problem{2-4}{物体放在凸球面反射镜前何处，可产生大小与物体相等的倒立实像？}
\answerarea
\clearpage
\problem{2-5}{凹面镜的半径为 $40\,\mathrm{cm}$，物体放在何处成放大两倍的实像？放在何处成放大两倍的虚像？}
\answerarea
\problem{2-6}{要把球面反射镜前 $10\,\mathrm{cm}$ 处的灯丝成像于 $3\,\mathrm m$ 处的墙上，镜形应是凸的还是凹的？半径应有多大？这时像放大了多少倍？}
\answerarea
\clearpage
\problem{2-7}{按已约定的正负号法则（I）、（II）、（III）、（IV）、（V）等，标示出本题各图中的物距 $s$、像距 $s'$、曲率半径 $r$、光线倾角 $u$、$u'$ 的绝对值。比较各图中折射率 $n$、$n'$ 的大小，指明各图中物像的虚实。\fig[0.85]{2_7}}
\answerarea
\problem{2-8}{分别根据上题各图推导球面折射成像公式（2.19）。}
\answerarea
\clearpage
\problem{2-9}{若空气中一均匀球形透明体能将平行光束会聚于其背面的顶点上，此透明体的折射率应等于多少？}
\answerarea
\problemfig[0.31]{2-10}{如本题图，一平行平面玻璃板的折射率为 $n$，厚度为 $h$，点光源 $Q$ 发出的傍轴光束（即接近于正入射的光束）经上表面反射，成像于 $Q'_1$；而折射线穿过上表面后在下表面反射，再从上表面折射的光束成像于 $Q'_2$。证明 $Q'_1$ 与 $Q'_2$ 间的距离为 $2h/n$。
\par【提示：把平面看成 $r\to\infty$ 的球面，并利用球面折射公式计算。】}{2_10}
\answerarea
\clearpage
\problemfig[0.43]{2-11}{如本题图，一会聚光束本来交于 $P$ 点，插入一折射率为 $1.50$ 的平面平行玻璃板后，像点移至 $P'$。求玻璃板的厚度。}{2_11}
\answerarea
\problem{2-12}{根据费马原理推导傍轴条件下球面反射成像公式（2.23）。}
\answerarea
\clearpage
\problem{2-13}{根据费马原理推导傍轴条件下球面折射成像公式（2.19）。}
\answerarea
\clearpage
\problem{2-14}{求联立方程式（2.15）和（2.16）的解，证明它所代表的共轭点就是第1.4节中给出的齐明点。
\begin{empheq}[left=\empheqlbrace]{align}
\frac{s^2}{n^2(s+r)^2}-\frac{s'^2}{n'^2(s'-r)^2}&=0,\tag{2.15}\\[5pt]
\frac{1}{n^2(s+r)}-\frac{1}{n'^2(s'-r)}&=0.\tag{2.16}
\end{empheq}
其中 $s,s'$ 为依原书正负号约定的物距与像距，$r$ 为球面曲率半径，$n,n'$ 分别为物方和像方介质的折射率。
\par\smallskip\textbf{齐明点：}图2-6中的 $Q,Q'$ 这对共轭点称为折射球面的齐明点（aplanatic points）。图中
\[u=\angle MQA,\qquad u'=\angle MQ'A,\qquad
\frac{\sin u}{\sin u'}=\frac{\overline{AQ'}}{\overline{QA}}\quad\text{（2.4）}.\]
\bookfig[0.38]{book_2_6}{图2-6\quad 折射球面的齐明点}}
\answerarea
\clearpage
\problem{2-15}{某透镜用 $n=1.50$ 的玻璃制成，它在空气中的焦距为 $10.0\,\mathrm{cm}$，它在水中的焦距为多少？（水的折射率为 $4/3$。）}
\answerarea
\problem{2-16}{一薄透镜折射率为 $1.50$，光焦度为 $5.00\,\mathrm D$。将此镜浸入某液体，光焦度变为 $-1.00\,\mathrm D$，求此液体的折射率。}
\answerarea
\clearpage
\problemfig[0.27]{2-17}{用一曲率半径为 $20\,\mathrm{cm}$ 的球面玻璃和一平面玻璃黏合成空气透镜，将其浸于水中（见本题图）。设玻璃壁厚可忽略，水和空气的折射率分别为 $4/3$ 和1，求此透镜的焦距。此透镜是会聚的还是发散的？}{2_17}
\answerarea
\problem{2-18}{一凸透镜的焦距为 $12\,\mathrm{cm}$，填充下表中的空白，并作出相应的光路图。\lenstable}
\answerarea
\clearpage
\problem{2-19}{一凹透镜的焦距为 $12\,\mathrm{cm}$，填充下表中的空白，并作出相应的光路图。\lenstable}
\answerarea
\problem{2-20}{在 $5\,\mathrm{cm}$ 焦距的凸透镜前放一小物，要虚像成像于 $25\,\mathrm{cm}$ 到无穷远之间，物应放在什么范围？}
\answerarea
\clearpage
\problem{2-21}{一光源与屏间的距离为 $1.6\,\mathrm m$，用焦距为 $30\,\mathrm{cm}$ 的凸透镜插在二者之间，透镜应放在什么位置，才能使光源成像于屏上？}
\answerarea
\problem{2-22}{屏幕放在距物 $100\,\mathrm{cm}$ 远处，二者之间放一凸透镜。当前后移动透镜时，我们发现透镜有两个位置可以使物成像在屏幕上。测得这两个位置之间的距离为 $20.0\,\mathrm{cm}$，求：（1）这两个位置到幕的距离和透镜的焦距；（2）两个像的横向放大率。}
\answerarea
\clearpage
\problem{2-23}{如上题，在固定的物与幕之间移动凸透镜。证明：要使透镜有两个成像位置，物和幕之间的距离必须大于四倍焦距。}
\answerarea
\problemfig[0.41]{2-24}{如本题图，$L_1$、$L_2$ 分别为凸透镜和凹透镜。前放一小物，移动屏幕到 $L_2$ 后 $20\,\mathrm{cm}$ 的 $S_1$ 处接到像。现将凹透镜 $L_2$ 撤去，将屏移前 $5\,\mathrm{cm}$ 至 $S_2$ 处，重新接收到像。求凹透镜 $L_2$ 的焦距。}{2_24}
\answerarea
\clearpage
\problem{2-25}{一光学系统由一焦距为 $5.0\,\mathrm{cm}$ 的会聚透镜 $L_1$ 和一焦距为 $10.0\,\mathrm{cm}$ 的发散透镜 $L_2$ 组成，$L_2$ 在 $L_1$ 之右 $5.0\,\mathrm{cm}$。在 $L_1$ 之左 $10.0\,\mathrm{cm}$ 处放置一小物，求经此光学系统后所成的像的位置和横向放大率。用作图法验证计算结果。}
\answerarea
\problem{2-26}{用逐次成像法解第4.4节中的例题4，并将结果与之比较。
\bookfig[0.94]{book_2_29}{图2-29\quad 第4.4节例题4配图}
\missing{例题4的完整题干与已知参数；图2-29已补入。}}
\answerarea
\clearpage
\problem{2-27}{用作图法求本题各图中 $PQ$ 的像（入射线从左到右）。\fig[0.9]{2_27}}
\answerarea
\problem{2-28}{用作图法求本题图中 $PQ$ 的像（入射线从左到右）。\fig[0.7]{2_28}}
\answerarea
\clearpage
\problem{2-29}{用作图法求本题图中联合光具组的主面和焦点。\fig[0.75]{2_29}}
\answerarea
\problem{2-30}{用作图法求本题图中正方形 $ABCD$ 的像。\fig[0.65]{2_30}}
\answerarea
\clearpage
\problem{2-31}{验算第7.7节图2-63a和b中给出的两种拉姆斯登目镜的主面与焦点。
\par\smallskip 向场镜 $L_1$ 和接目镜 $L_2$ 的焦距和间隔之比，在图（a）中为
\[f'_1:f'_2:d=1:1:1.\]
图中 $F,F'$ 为整个系统的焦点，$H,H'$ 为主点，$\mathcal H,\mathcal H'$ 表示相应的主平面；物体 $QP$ 放在焦点 $F$ 内侧附近。图（b）为缩短两镜间距的结构示意。
\bookfig[0.98]{book_2_63}{图2-63\quad 拉姆斯登目镜}}
\answerarea
\problem{2-32}{求下表中厚透镜的焦距和主面、焦点的位置，并作图表示。已知玻璃的折射率为 $1.500$，两界面顶点间的距离为 $1.00\,\mathrm{cm}$，透镜放在空气中。
\begin{center}\begin{tabular}{c|c|r|r}\hline &形状&$r_1/\mathrm{cm}$&$r_2/\mathrm{cm}$\\\hline
（1）&双凸&10.0&$-10.0$\\（2）&凸凹&10.0&20.0\\（3）&凹凸&$-15.0$&$-20.0$\\\hline\end{tabular}\end{center}}
\answerarea
\clearpage
\problem{2-33}{求放在空气中玻璃球的焦距和主面、焦点的位置，并作图表示。已知玻璃球的半径为 $2.00\,\mathrm{cm}$，折射率为 $1.500$。}
\answerarea
\problem{2-34}{上题中玻璃球表面上有一斑点，计算从另一侧观察此斑点像的位置和放大率，并用作图法验证之。}
\answerarea
\clearpage
\problemfig[0.42]{2-35}{（1）用作图法求本题图中光线1的共轭线；（2）在图上标出光具组节点 $N$、$N'$ 的位置。}{2_35}
\answerarea
\problemfig[0.42]{2-36}{已知 $1-1'$ 是一对共轭光线，求光线2的共轭线（见本题图）。}{2_36}
\answerarea
\clearpage
\problem{2-37}{对一光具组，测得当物距改变 $\Delta x$ 时，像距改变 $\Delta x'$，同时横向放大率由 $V_1$ 变为 $V_2$。试证明此光具组的焦距为
\[f=\frac{\Delta x}{\dfrac1{V_2}-\dfrac1{V_1}},\qquad f'=\frac{\Delta x'}{V_2-V_1}.\]
这里提供了一种测焦距的方法，它与测焦点位置的方法配合起来，可以确定光具组的主面。}
\answerarea
\problem{2-38}{一架幻灯机的投影镜头的焦距为 $7.5\,\mathrm{cm}$，当幕由 $8\,\mathrm m$ 移至 $10\,\mathrm m$ 远时，镜头移动多少距离？}
\answerarea
\clearpage
\problem{2-39}{某照相机与拍摄物体的最近距离为 $1\,\mathrm m$，装上光焦度为 $2\,\mathrm D$ 的近视镜片，能拍摄的最近距离为多少？（假设近拍镜与照相机镜头是密接的。）}
\answerarea
\problem{2-40}{某人对 $2.5\,\mathrm m$ 以外的物看不清，需配多少度的眼镜？另一人对 $1\,\mathrm m$ 以内的物看不清，需配怎样的眼镜？}
\answerarea
\clearpage
\problem{2-41}{计算 $2\times$、$3\times$、$5\times$、$10\times$ 放大镜或目镜的焦距。}
\answerarea
\problem{2-42}{一架显微镜，物镜焦距为 $4\,\mathrm{mm}$，中间像放在物镜像方焦点后面 $160\,\mathrm{mm}$ 处，如果目镜是 $20\times$ 的，显微镜的总放大率是多少？}
\answerarea
\clearpage
\problem{2-43}{一架显微镜的物镜和目镜相距 $20.0\,\mathrm{cm}$，物镜焦距为 $7.0\,\mathrm{mm}$，目镜焦距为 $5.0\,\mathrm{mm}$，把物镜和目镜都看成是薄透镜。（1）求被观测物到物镜的距离；（2）物镜的横向放大率；（3）显微镜的总放大率；（4）焦深。}
\answerarea
\problem{2-44}{物镜、目镜皆为会聚的望远镜称为开普勒型望远镜（图2-29），物镜会聚而目镜发散的望远镜，称为伽里略型望远镜。
\begin{enumerate}\item 画出伽里略型望远镜的光路。\item 一伽里略型望远镜的物镜和目镜相距 $12\,\mathrm{cm}$，若望远镜的放大率为 $4\times$，物镜和目镜的焦距各多少？\end{enumerate}}
\answerarea
\clearpage
\problem{2-45}{拟制一个 $3\times$ 的望远镜，已有一个焦距为 $50\,\mathrm{cm}$ 的物镜，问，在（1）开普勒型和（2）伽里略型中，目镜的光焦度以及物镜、目镜的距离各多少？}
\answerarea
\problem{2-46}{一望远镜的物镜直径为 $5.0\,\mathrm{cm}$，焦距为 $20\,\mathrm{cm}$，目镜直径为 $1.0\,\mathrm{cm}$，焦距为 $2.0\,\mathrm{cm}$，求此望远镜的入射光瞳和出射光瞳的位置和大小。}
\answerarea
\clearpage
\problem{2-47}{望远镜的孔径光阑和入射光瞳通常就是其物镜的边缘。求出射光瞳的位置，并证明出射光瞳直径 $D'$ 与物镜直径 $D_0$ 之比为
\[D'=\frac{D_0}{|M|},\]
其中 $M=-f'_{\mathrm{ob}}/f'_{\mathrm{oc}}$ 是望远镜的视角放大率。}
\answerarea
\problem{2-48}{将望远镜倒过来可作激光扩束之用。设一望远镜物镜焦距为 $30\,\mathrm{cm}$，目镜焦距为 $1.5\,\mathrm{cm}$，它能使激光光束的直径扩大几倍？}
\answerarea
\clearpage
\problem{2-49}{显微镜的孔径光阑和入射光瞳通常就是其物镜的边缘。求：（1）出射光瞳的位置；（2）证明在傍轴近似下出射光瞳的直径 $D'$ 与入射孔径角 $u_0$ 的关系是
\[D'=\frac{2s_0n u_0}{|M|},\]
式中 $s_0=25\,\mathrm{cm}$ 是明视距离，$M$ 是显微镜的视角放大率，$n$ 是物方折射率（除油浸物镜外，$n\approx1$）。严格的公式应把 $u_0$ 换成 $\sin u_0$，参见习题2-56。\missing{原截图脚注的“2-56”疑指本批2-53，保留原引用待核。}}
\answerarea
\problem{2-50}{如本题图中 $L_1$、$L_2$ 是两个会聚透镜，$Q$ 是物点，$DD$ 是光阑。已知焦距 $f_1=2a$，$f_2=a$，图中标示各距离为 $s=10a$、$l=4a$、$d=6a$；此外透镜与光阑半径之比是 $r_1=r_2=3r_3$。求此光具组的孔径光阑、入射光瞳、出射光瞳、入射窗和视场光阑的位置和大小。\fig[0.9]{2_50}}
\answerarea
\clearpage
\problemfig[0.39]{2-51}{惠更斯目镜的结构详见第7.7节图2-62。如本题图所示，今在两透镜间放一光阑 $AA$，设透镜 $L_1$、$L_2$ 和光阑的直径分别为 $D_1$、$D_2$ 和 $D$。试证：
\begin{enumerate}\item 向场镜 $L_1$ 成为孔径光阑的条件为 $D_1<D_2$；\item 光阑 $AA$ 成为视场光阑的条件是 $D/2<D_2$；\item 这时对 $F'_2$ 点计算的出射孔径角 $u'_0$ 由下式确定：$\tan u'_0=D_1/(2a)$。\end{enumerate}}{2_51}
\answerarea
\problem{2-52}{证明在一对齐明点附近不可能有另一对齐明点。}
\answerarea
\clearpage
\problem{2-53}{试证明，显微镜出射光瞳直径 $D'$ 由下式确定：
\[D'=\frac{2s_0n\sin u_0}{|M|}=2s_0\frac{\mathrm{N.A.}}{|M|},\]
即 $D'$ 正比于数值孔径 $\mathrm{N.A.}=n\sin u_0$，反比于视角放大率 $M$。（$s_0=25\,\mathrm{cm}$ 是明视距离。）}
\answerarea
\problem{2-54}{拟用冕牌玻璃 K9（$n_D=1.5163$，$n_F=1.5220$，$n_C=1.5139$）和重火石玻璃 F4（$n_D=1.6199$，$n_F=1.6321$，$n_C=1.6150$）来做消色差胶合透镜，焦距为 $100\,\mathrm{mm}$。若已确定其负透镜的非黏合面为平面，试求其余各面的曲率半径。}
\answerarea
\clearpage
\problem{2-55}{太阳表面的辐射亮度为 $2\times10^7\,\mathrm{W/(m^2\cdot sr)}$，用相对孔径 $D/f=1.5$ 的放大镜将阳光聚焦成光斑的最大辐射照度为多少？}
\answerarea
\problem{2-56}{一屏放在离烛 $100\,\mathrm{cm}$ 处，把一会聚薄透镜放在烛和屏之间，透镜有两个位置可以在屏上得到烛的像，这两个位置相距 $20\,\mathrm{cm}$，在屏上烛的两个像的照度相差多少倍？}
\answerarea
\clearpage
\problem{2-57}{望远镜物镜的直径为 $75\,\mathrm{mm}$，当放大率为（1）20倍，（2）25倍，（3）50倍时，求望远镜中月亮的像的主观亮度与天然主观亮度之比。眼瞳瞳孔的直径为 $3.0\,\mathrm{mm}$。}
\answerarea
\problem{2-58}{一天文望远镜的物镜直径等于 $18\,\mathrm{cm}$，透光系数为 $0.50$。已知肉眼可直接观察到六等星。求：
\begin{enumerate}\item 用此望远镜所能看到的最弱星等；\item 最适宜观察星的放大率（正常放大率）；\item 当放大率为10倍时可以见到星的等次。\end{enumerate}
设眼睛瞳孔的直径可取 $3.0\,\mathrm{mm}$。【注：星等增加一等，其亮度减小 $\sqrt[5]{100}\approx2.5$ 倍。】}
\answerarea
\clearpage
\problem{2-59}{求数值孔径为 $1.5$ 的显微镜的正常放大率，设瞳孔直径为 $3.0\,\mathrm{mm}$。}
\answerarea
\clearpage

\chapter{干涉}
\problem{3-1}{在杨氏双孔实验中，孔距为 $0.1\,\mathrm{mm}$，孔与屏幕的距离为 $3\,\mathrm m$，对下列三条典型谱线求出干涉条纹的间距：F蓝线（$486.1\,\mathrm{nm}$），D黄线（$589.3\,\mathrm{nm}$），C红线（$656.3\,\mathrm{nm}$）。}
\answerarea
\problem{3-2}{在杨氏双孔实验中，孔距为 $0.45\,\mathrm{mm}$，孔与幕的距离为 $1.2\,\mathrm m$，测得10个亮纹之间的间距为 $1.5\,\mathrm{cm}$，问光源的波长是多少。}
\answerarea
\clearpage
\problem{3-3}{一微波检测器安装在湖滨高出水面 $0.5\,\mathrm m$ 处。当一颗发射 $21\,\mathrm{cm}$ 波长单色微波的射电星体徐徐自地平线升起时，检测器指出一系列信号强度的极大和极小。当第一个极大出现时，射电星体相对地平线的仰角 $\theta$ 为多少？}
\answerarea
\problemfig[0.4]{3-4}{本题图所示装置是昆克（G. Quincke）用来测声波波长的。管口 $T$ 置于单一声源的声源之前，声波分 $A$、$B$ 两股传播到出口 $O$，其中一股 $B$ 的长度像乐队中的长号那样，可以拉出拉进。当 $A$、$B$ 两股等长时，声音近似保持原有的强度。当 $A$ 逐渐拉开到 $d=16.0\,\mathrm{cm}$ 时，在管口 $O$ 处的声音第一次消失，求此声波的波长。}{3_4}
\answerarea
\clearpage
\problem{3-5}{设菲涅耳双镜的夹角为 $20'$，缝光源离双镜交线 $1.0\,\mathrm{cm}$，接收屏幕与光源经双镜所成的两个虚像联线平行，幕与两镜交线的距离为 $210\,\mathrm{cm}$，光波长为 $600.0\,\mathrm{nm}$。问：
\begin{enumerate}\item 干涉条纹的间距为多少？\item 在幕上最多能看到几根干涉条纹？\item 如果光源到两镜交线的距离增大一倍，干涉条纹有什么变化？\item 如果光源与两镜交线的距离保持不变，而在横向有所移动，干涉条纹有什么变化？\item 如果要在幕上出现有一定衬比度的干涉条纹，允许缝光源的最大宽度为多少？\end{enumerate}
}
\answerarea
\clearpage
\problemfig[0.42]{3-6}{点光源置于薄透镜的焦点，薄透镜后放一双棱镜（见本题图）。设双棱镜的顶角为 $3'30''$，折射率为 $1.5$，屏幕与棱镜相距 $5.0\,\mathrm m$，光波长为 $500.0\,\mathrm{nm}$。求幕上条纹的间距。幕上能出现几根干涉条纹？}{3_6}
\answerarea
\problem{3-7}{设劳埃镜的镜长有 $5.0\,\mathrm{cm}$，幕与镜边缘的距离为 $3.0\,\mathrm m$，缝光源离镜面高度为 $0.5\,\mathrm{mm}$，水平距离 $2.0\,\mathrm{cm}$，光波长为 $589.3\,\mathrm{nm}$。求幕上条纹的间距。幕上能出现几根干涉条纹？}
\answerarea
\clearpage
\problem{3-8}{本题图为梅斯林（L. Meslin）干涉装置，将透镜对剖后再沿光轴方向将两半 $L_1$、$L_2$ 错开一定距离。光点 $S$ 位于光轴上，$S'_1$、$S'_2$ 是它的像。
\begin{enumerate}\item 在图上标出相干光束的交叠区；\item 在交叠区中放一屏幕垂直于光轴，幕上干涉条纹的形状是怎样的？\item 设透镜焦距为 $30\,\mathrm{cm}$，$S$ 与 $L_1$ 的距离为 $60\,\mathrm{cm}$，$L_1$ 与 $L_2$ 的距离为 $8.0\,\mathrm{cm}$，光波长为 $500.0\,\mathrm{nm}$。两像间的中点距离透镜 $L_2$ 有多远？在此处放一屏幕，在其上接收到的亮纹间距为多少？\end{enumerate}\fig[0.65]{3_8}}
\answerarea
\clearpage
\problemfig[0.35]{3-9}{本题图所示为一种利用干涉现象测定气体折射率的原理性结构。在 $S_1$ 孔后面放置一长度为 $l$ 的透明容器，当被测气体注入容器而将空气赶出的过程中幕上的干涉条纹就会移动。由移过条纹的根数即可推知气体的折射率。
\begin{enumerate}\item 设待测气体的折射率大于空气折射率，干涉条纹如何移动？\item 设 $l=2.0\,\mathrm{cm}$，条纹移过20根，光波长为 $589.3\,\mathrm{nm}$，空气折射率为 $1.000276$，求待测气体（氯气）的折射率。\end{enumerate}}{3_9}
\answerarea
\problemfig[0.42]{3-10}{瑞利（Rayleigh）干涉仪的结构和使用原理如下（参见本题图）：以钠光灯作为光源置于透镜 $L_1$ 的前焦点，在透镜 $L_2$ 的后焦面上观测干涉条纹的变动。在两个透镜之间安置一对完全相同的玻璃管 $T_1$ 和 $T_2$。实验开始时，$T_2$ 管充以空气，$T_1$ 管抽成真空，此时开始观测干涉条纹。然后逐渐使空气进入 $T_1$ 管，直到它与 $T_2$ 管的气压相同为止。记下这一过程中条纹移动的数目。设光波长为 $589.3\,\mathrm{nm}$，管长 $20\,\mathrm m$，条纹移动98根，求空气的折射率。}{3_10}
\answerarea
\clearpage
\problem{3-11}{用钠光灯作杨氏双缝干涉实验。光源宽度被限制为 $2\,\mathrm{mm}$，带双缝的幕离缝光源 $2.5\,\mathrm m$，为了在幕上获得可见的干涉条纹，双缝间隔不能大于多少？}
\answerarea
\problem{3-12}{一个直径为 $1\,\mathrm{cm}$ 的发光面元，如果用干涉孔径角量度的话，其空间相干性是多少弧度？如果用相干面积量度，则 $1\,\mathrm m$ 远的相干面积为多大？$10\,\mathrm m$ 远的相干面积为多大？}
\answerarea
\clearpage
\problemfig[0.29]{3-13}{把直径为 $D$ 的细丝夹在两块平玻璃的一边，形成尖劈形空气层（见本题图下方）。在钠黄光（$\lambda=589.3\,\mathrm{nm}$）的垂直照射下形成如图上方所示的干涉条纹。试问 $D$ 为多少？}{3_13}
\answerarea
\clearpage
\problemfig[0.38]{3-14}{块规是机械工里用的一种长度标准，它是一钢质长方体，它的两个端面经过精密抛光，达到相互平行。本题图中 $G_1$、$G_2$ 是同规格的两个块规，$G_1$ 的长度是标准的，$G_2$ 是要校准的。校准的方法如下：把 $G_1$ 和 $G_2$ 放在钢质平台面上，使侧面和底部接触。$G_1$、$G_2$ 上面用一块透明平板 $T$ 压住。如果 $G_1$ 和 $G_2$ 的高度（即长度）不等，略有差别，则在 $T$ 和 $G_1$、$G_2$ 之间分别形成尖劈形空气层。它们在单色光照射下产生等厚干涉条纹。
\begin{enumerate}\item 设入射光的波长是 $589.3\,\mathrm{nm}$，$G_1$ 和 $G_2$ 相隔 $5\,\mathrm{cm}$（即图中的 $l$），$T$ 和 $G_1$、$G_2$ 间干涉条纹的间隔都是 $0.5\,\mathrm{mm}$。试求块规 $G_2$ 和 $G_1$ 的高度之差。怎样判断它们谁高谁矮？\item 如果 $T$ 和 $G_1$ 间干涉条纹的间隔是 $0.5\,\mathrm{mm}$，而 $T$ 和 $G_2$ 间的是 $0.3\,\mathrm{mm}$，则说明什么问题？\end{enumerate}}{3_14}
\answerarea
\clearpage
\problem{3-15}{在图3-27所示的干涉膨胀计中，样品与石英环的高度均为 $1\,\mathrm{cm}$，当温度升高 $100^\circ\mathrm C$ 时，视场中的干涉条纹移动20根。求样品的线膨胀系数。设光波长为 $589.3\,\mathrm{nm}$，石英的线膨胀系数为 $0.35\times10^{-6}/{}^\circ\mathrm C$。
\bookfig[0.43]{book_3_27}{图3-27\quad 干涉膨胀计}
\par{\small 图中 $G$ 为标准石英环，$C$ 为待测柱形样品，$\Pi_1$ 为上方石英平板，$\Pi_2$ 为下方支承平板。温度改变时，样品上端与 $\Pi_1$ 之间楔形空气层的厚度改变，可通过干涉条纹移动测量样品与石英环长度的相对变化。}}
\answerarea
\clearpage
\problem{3-16}{本题图a所示为一种测 PN 结的结深 $x_j$ 的方法。在 N 型半导体基质硅片表面经杂质扩散而形成 P 型半导体区。P 区与 N 区的交界面叫 PN 结，PN 结距表面的深度（即 P 区的厚度 $x_j$）叫做结深，在半导体工艺上需要测定结深。测量的方法是先通过磨角、染色，使 P 区和 N 区的分界线稍微地显示出来，然后盖上半反射膜，在它与硅片之间形成尖劈形空气薄膜。用单色光垂直照射时，可以观察到空气薄膜的等厚干涉条纹。数出 P 区空气薄膜的条纹数目 $\Delta k$ 即可求出结深
\[x_j=\Delta k\frac{\lambda}{2}.\]
由于光在金属或半导体表面反射时相位变化比较复杂，用本方法测量结深 $x_j$ 没有考虑此相位突变，因此测量结果不太精确。更精确的测量方法见图b，半反射膜不是盖在 P 区的上面，而是稍微往上翘一点，观察到的干涉条纹如图b下方所示。试说明：（1）干涉条纹为什么会是这样的？（2）若用 $\lambda=550.0\,\mathrm{nm}$ 的光测得斜干涉条纹的间隔为 $0.20\,\mathrm{mm}$，交界线上两点 $AB$ 间的距离为 $1.1\,\mathrm{mm}$，结深 $x_j$ 为多少？（3）此法比图a所示的方法精确在哪里？\fig[0.78]{3_16}}
\answerarea
\clearpage
\problem{3-17}{测得牛顿环从中间数第5环和第15环的半径分别为 $0.70\,\mathrm{mm}$ 和 $1.7\,\mathrm{mm}$。求透镜的曲率半径。设光波长为 $0.83\,\mu\mathrm m$。}
\answerarea
\problem{3-18}{肥皂膜的反射光呈现绿色，这时膜的法线和视线的角约为 $35^\circ$，试估算膜的最小厚度。设肥皂水的折射率为 $1.33$，绿光波长为 $500.0\,\mathrm{nm}$。}
\answerarea
\clearpage
\problem{3-19}{在玻璃表面上涂一层折射率为 $1.30$ 的透明薄膜，设玻璃折射率为 $1.5$。
\begin{enumerate}\item 对于波长为 $500.0\,\mathrm{nm}$ 的入射光来说，膜厚应为多少才能使反射光干涉相消？这时强度反射率为多少？\item 对波长为 $400.0\,\mathrm{nm}$ 的紫光和 $700.0\,\mathrm{nm}$ 的红光来说，第（1）问所得的厚度在两束反射相干光之间产生多大的相位差？（不考虑色散。）\end{enumerate}
\editorial{原材料第（2）问交叉引用印作“第（2）问”，此处按题意改为“第（1）问”。}}
\answerarea
\clearpage
\problem{3-20}{砷化镓发光管制成半球形，以增加位于球心的发光区对外输出功率，减少反射损耗（见本题图）。已知砷化镓发射光波长为 $930\,\mathrm{nm}$，折射率为 $3.4$。为了进一步提高输出功率，常在球形表面涂敷一层增透膜。
\begin{enumerate}\item 不加增透膜时，球面的强度反射率有多大？\item 增透膜的折射率和厚度应取多大？\item 如果用氟化镁（折射率为 $1.38$）能否增透？强度反射率有多大？\item 如果用硫化锌（折射率为 $2.35$）能否增透？强度反射率有多大？\end{enumerate}\fig{3_20}}
\answerarea
\clearpage
\problem{3-21}{如图3-33用肉眼直接观察薄膜表面的干涉条纹。设瞳孔直径为 $3\,\mathrm{mm}$，与表面相距 $30\,\mathrm{cm}$，视线与表面法线夹角 $30^\circ$，薄膜折射率为 $1.5$。
\begin{enumerate}\item 分别计算膜厚 $2\,\mathrm{cm}$ 及 $20\,\mu\mathrm m$ 两种情况下，点源 $Q_1$、$Q_2$ 在观察点 $P$ 产生的光程差改变量 $\delta(\Delta L)$。\item 如果为了保证条纹有一定的衬比度，要求上述光程差改变量的数量级不能超过多少？以此来估计一下对膜厚 $h$ 的限制。\end{enumerate}
\bookfig[0.48]{book_3_33}{图3-33\quad 干涉条纹衬比度下降的原因}}
\answerarea
\problem{3-22}{证明迈克耳孙干涉仪中圆形等倾条纹的半径与整数的平方根成正比。}
\answerarea
\clearpage
\problem{3-23}{用钠光（$589.3\,\mathrm{nm}$）观察迈克耳孙干涉条纹，先看到干涉场中有12个亮环，且中心是亮的。移动平面镜 $M_1$ 后，看到中心吞（吐）了10环，而此时干涉场中还剩有5个亮环。试求：（1）$M_1$ 移动的距离；（2）开始时中心亮斑的干涉级；（3）$M_1$ 移动后，从中心向外数第5个亮环的干涉级。}
\answerarea
\problem{3-24}{在迈克耳孙干涉仪中，反射镜移动 $0.33\,\mathrm{mm}$，测得条纹变动192次，求光的波长。}
\answerarea
\clearpage
\problem{3-25}{钠光灯发射的黄线包含两条相近的谱线，平均波长为 $589.3\,\mathrm{nm}$。在钠光下调节迈克耳孙干涉仪，人们发现干涉场的衬比度随镜面移动而周期性地变化。实测的结果由条纹最清晰到最模糊，视场中吞（吐）490圈条纹，求钠双线的两个波长。}
\answerarea
\problem{3-26}{在一次迈克耳孙干涉仪实验中，所用的线纹标准具长度为 $0.39\,\mathrm{mm}$，如用镉灯（$643.847\,\mathrm{nm}$）作光源，实验时所测得的条纹变动数目应是多少？}
\answerarea
\clearpage
\problem{3-27}{用迈克耳孙干涉仪进行精密测长，光源为 $632.8\,\mathrm{nm}$ 的氦氖激光，其谱线宽度为 $10^{-4}\,\mathrm{nm}$，整机接收（光电转换）灵敏度可达 $1/10$ 个条纹。求这台仪器测长精度为多少？一次测长量程为多少？}
\answerarea
\problem{3-28}{迈克耳孙干涉仪中的一臂（反射镜）以速度 $v$ 匀速推移，用透镜接收干涉条纹，将它会聚到光电元件上，把光强变化转换为电信号。
\begin{enumerate}\item 若测得电信号的时间频率为 $\nu$，求入射光的波长 $\lambda$。\item 若入射光波长在 $0.8\,\mu\mathrm m$ 左右，要使电信号频率控制在 $50\,\mathrm{Hz}$，反射镜平移的速度应为多少？\item 按以上速度移动反射镜，钠黄光产生电信号的拍频为多少？（钠黄光双线波长为 $589.0\,\mathrm{nm}$ 和 $589.6\,\mathrm{nm}$。）\end{enumerate}}
\answerarea
\clearpage
\problem{3-29}{有两个波长 $\lambda_1$ 和 $\lambda_2$，在 $600.0\,\mathrm{nm}$ 附近相差 $0.0001\,\mathrm{nm}$，要用法布里--珀罗干涉仪把它们分辨开来，间隔 $h$ 需要多大？设反射率 $R=0.95$。}
\answerarea
\problem{3-30}{如果法布里--珀罗干涉仪两反射面之间的距离为 $1.0\,\mathrm{cm}$，用绿光（$500.0\,\mathrm{nm}$）做实验，干涉图样的中心正好是一亮斑。求第10个亮环的角直径。}
\answerarea
\clearpage
\problem{3-31}{设法--珀腔长 $5\,\mathrm{cm}$，用扩展光源做实验，光波波长为 $0.6\,\mu\mathrm m$，问：
\begin{enumerate}\item 中心干涉级数为多少？\item 在倾角为 $1^\circ$ 附近干涉环的半角宽度为多少？设反射率 $R=0.98$。\item 如果用这个法--珀腔分辨谱线，其色分辨本领有多高？可分辨的最小波长间隔有多少？\item 如果用这个法--珀腔对日光进行选频，透射最强的谱线有几条？每条谱线宽度有多少？\item 由于热胀冷缩，引起腔长的改变量为 $10^{-5}$（相对值），求谱线的漂移量（相对值）。\end{enumerate}}
\answerarea
\clearpage
\problemfig[0.39]{3-32}{利用多光束干涉可以制成一种干涉滤光片。如本题图，在很平的玻璃片上镀一层银，在银面上加一层透明膜，例如冰晶石（$3\mathrm{NaF}\cdot\mathrm{AlF}_3$），其上再镀一层银。设银面的反射率 $R=0.96$，于是两个银面之间就形成一个膜层，产生多光束干涉。透明膜的折射率为 $1.55$，膜厚 $h=4\times10^{-5}\,\mathrm{cm}$，平行光正入射。问：（1）在可见光范围内，透射最强的谱线有几条？（2）每条谱线宽度为多少？}{3_32}
\answerarea
\clearpage

\chapter{衍射}
\problem{4-1}{在菲涅耳圆孔衍射实验中，圆孔半径 $2.0\,\mathrm{mm}$，光源离圆孔 $2.0\,\mathrm m$，波长 $0.5\,\mu\mathrm m$，当接收屏幕由很远的地方向圆孔靠近时，求：（1）前三次出现中心亮斑（强度极大）的位置；（2）前三次出现中心暗斑（强度极小）的位置。}
\answerarea
\problem{4-2}{在菲涅耳圆孔衍射实验中，光源距离圆孔 $1.5\,\mathrm m$，波长 $0.63\,\mu\mathrm m$，接收屏幕与圆孔距离 $6.0\,\mathrm m$，圆孔半径从 $0.5\,\mathrm{mm}$ 开始逐渐扩大。求：（1）最先的两次出现中心亮斑时圆孔的半径；（2）最先的两次出现中心暗斑时圆孔的半径。}
\answerarea
\clearpage
\problem{4-3}{用直刀口将点光源的波面遮住一半（直边衍射），几何阴影边缘点上的光强比自由传播时小多少倍？}
\answerarea
\problem{4-4}{求圆孔中露出 $1.5$ 个半波带时衍射场中心强度与自由传播时强度之比。}
\answerarea
\clearpage
\problem{4-5}{用平行光照明衍射屏，如对波面作如本题图所示的下几种方式的遮挡，求轴上点的光强与自由传播时之比（图中标出的是该处到场点的光程，其中 $b$ 是中心到场点的光程）。\fig[0.85]{4_5}}
\answerarea
\problem{4-6}{若一个菲涅耳波带片只将前五个偶数半波带遮挡，其余地方都开放，求衍射场中心强度与自由传播时之比。}
\answerarea
\clearpage
\problem{4-7}{若一个菲涅耳波带片将前50个奇数半波带遮挡，其余地方都开放，求衍射场中心强度与自由传播时之比。}
\answerarea
\problem{4-8}{菲涅耳波带片第一个半波带的半径 $\rho_1=5.0\,\mathrm{mm}$。（1）用波长 $\lambda=1.06\,\mu\mathrm m$ 的单色平行光照明，求主焦距；（2）若要求主焦距为 $25\,\mathrm{cm}$，需将此波带片缩小多少倍？}
\answerarea
\clearpage
\problem{4-9}{如何制作一张满足以下要求的波带片：（1）它在 $400.0\,\mathrm{nm}$ 紫光照明下的主焦距为 $80\,\mathrm{cm}$；（2）主焦点光强是自由传播时的 $10^4$ 倍左右。}
\answerarea
\problem{4-10}{一菲涅耳波带片对 $900.0\,\mathrm{nm}$ 的红外光主焦距为 $30\,\mathrm{cm}$，改用 $632.8\,\mathrm{nm}$ 的氦氖激光照明，主焦距为多少？}
\answerarea
\clearpage
\problem{4-11}{如本题图，平行光以 $\theta_0$ 角斜入射在宽度为 $a$ 的单缝上。试证明：
\begin{enumerate}\item 夫琅禾费衍射的强度公式基本不变（忽略倾斜因子），即
\[I_\theta=I_0\left(\frac{\sin\alpha}{\alpha}\right)^2,\qquad \alpha=\frac{\pi a}{\lambda}(\sin\theta-\sin\theta_0),\]
式中 $I_0$ 为零级中心强度，只不过 $\alpha$ 的定义与正入射不同；\item 零级中心的位置在几何光学像点处；\item 零级斑半角宽度为 $\Delta\theta=\lambda/(a\cos\theta_0)$。\end{enumerate}\fig[0.9]{4_11}
\editorial{对照PDF，第（3）问分母同样印为 $a\sin\theta_0$；按第（1）问在 $\theta=\theta_0$ 附近展开，分母应为 $a\cos\theta_0$，此处据此校正。}}
\answerarea
\clearpage
\problem{4-12}{估算下列情形光束在介质界面上反射和折射时反射光束和折射光束的衍射发散角。设界面的线度为 $1\,\mathrm{cm}$，光波波长为 $0.6\,\mu\mathrm m$，折射率为 $1.5$：（1）平行光正入射；（2）入射角为 $75^\circ$；（3）入射角为 $89^\circ$（掠入射）。}
\answerarea
\problem{4-13}{试用巴比涅原理证明：互补的衍射屏产生的夫琅禾费衍射图样相同。}
\answerarea
\clearpage
\problem{4-14}{衍射细丝测径仪就是把单缝夫琅禾费衍射装置中的单缝用细丝代替。今测得零级衍射斑的宽度（两个一级暗纹间的距离）为 $1\,\mathrm{cm}$，求细丝的直径。已知光波波长 $0.63\,\mu\mathrm m$，透镜焦距 $50\,\mathrm{cm}$。}
\answerarea
\problem{4-15}{一对双星的角间隔为 $0.05''$。（1）需要多大口径的望远镜才能分辨它们？（2）此望远镜的角放大率应设计为多少才比较合理？}
\answerarea
\clearpage
\problem{4-16}{一台天文望远镜的口径为 $2.16\,\mathrm m$，由这一数据你能进一步获得关于它在光学性能方面的哪些知识？}
\answerarea
\problem{4-17}{一台显微镜，已知其 $\mathrm{N.A.}=1.32$，物镜焦距 $f'_o=1.91\,\mathrm{mm}$，目镜焦距 $f'_E=50\,\mathrm{mm}$。求：（1）最小分辨距离；（2）有效放大率；（3）光学筒长。}
\answerarea
\clearpage
\problem{4-18}{用一架照相机在离地面 $200\,\mathrm{km}$ 的高空拍摄地面上的物体，如果要求它能分辨地面上相距 $1\,\mathrm m$ 的两点，照相机的镜头至少要多大？设镜头的几何像差已很好地消除，感光波长为 $400.0\,\mathrm{nm}$。}
\answerarea
\problem{4-19}{已知地月距离约为 $3.8\times10^5\,\mathrm{km}$，用口径为 $1\,\mathrm m$ 的天文望远镜能分辨月球表面两点的最小距离是多少？}
\answerarea
\clearpage
\problem{4-20}{已知日地距离约为 $1.5\times10^8\,\mathrm{km}$，要求分辨太阳表面相距 $20\,\mathrm{km}$ 的两点，望远镜的口径至少需有多大？}
\answerarea
\problem{4-21}{用坐标纸绘制 $N=2$、$d=3a$ 的夫琅禾费衍射强度分布曲线，横坐标取 $\sin\theta$，至少画到第7级主极大，并计算第0级和第1级主极大与单缝主极大之比。}
\answerarea
\clearpage
\problem{4-22}{用坐标纸绘制 $N=6$、$d=1.5a$ 的夫琅禾费衍射强度分布曲线，横坐标取 $\sin\theta$，至少画到第4级主极大，并计算第4级主极大与单缝主极大之比。}
\answerarea
\problemfig[0.37]{4-23}{导出斜入射时夫琅禾费多缝衍射强度分布公式：
\[I_\theta=I_0\left(\frac{\sin\alpha'}{\alpha'}\right)^2\left(\frac{\sin N\beta'}{\sin\beta'}\right)^2,\]
其中 $\alpha'=\dfrac{\pi a}{\lambda}(\sin\theta-\sin\theta_0)$，$\beta'=\dfrac{\pi d}{\lambda}(\sin\theta-\sin\theta_0)$。$\theta_0$ 为入射线与光轴的夹角（见本题图）。}{4_23}
\answerarea
\clearpage
\problem{4-24}{导出斜入射时主极大位置公式、第 $k$ 级主极大的半角宽度公式及缺级情况，并注意与正入射情况作比较。}
\answerarea
\problem{4-25}{有三条平行狭缝，宽度都是 $a$，缝距分别为 $d$ 和 $2d$（见本题图）。证明正入射时夫琅禾费衍射强度分布公式为
\[I_\theta=I_0\left(\frac{\sin\alpha}{\alpha}\right)^2[3+2(\cos2\beta+\cos4\beta+\cos6\beta)],\]
其中 $\alpha=\pi a\sin\theta/\lambda$，$\beta=\pi d\sin\theta/\lambda$。\fig{4_25}}
\answerarea
\clearpage
\problem{4-26}{导出不等宽双缝的夫琅禾费衍射强度分布公式，缝宽分别为 $a$ 和 $2a$，缝距 $d=3a$（见本题图）。\fig{4_26}}
\answerarea
\problemfig[0.43]{4-27}{有 $2N$ 条平行狭缝，缝宽相同都是 $a$，缝间不透明部分的宽度作周期性变化 $a$、$3a$、$a$、$3a$、$\ldots$（见本题图）。求下列各种情形中的衍射强度分布：（1）遮住偶数缝；（2）遮住奇数缝；（3）全开放。}{4_27}
\answerarea
\clearpage
\problem{4-28}{波长为 $650.0\,\mathrm{nm}$ 的红光谱线，被观察发现它是双线。如果在 $9\times10^4$ 条刻线光栅的第3级光谱中恰好能分辨此双线，求其波长差。}
\answerarea
\problem{4-29}{若要50条/mm的光栅在第2级光谱中能分辨钠双线 $\lambda_1=589.0\,\mathrm{nm}$ 和 $\lambda_2=589.6\,\mathrm{nm}$，光栅宽度应选多少？}
\answerarea
\clearpage
\problem{4-30}{绿光 $500.0\,\mathrm{nm}$ 正入射在光栅常数为 $2.5\times10^{-4}\,\mathrm{cm}$、宽度为 $3\,\mathrm{cm}$ 的光栅上，聚光镜的焦距为 $50\,\mathrm{cm}$。（1）求第1级光谱的线色散；（2）求第1级光谱中能分辨的最小波长差；（3）设光栅最多能看到第几级光谱？}
\answerarea
\problem{4-31}{一束白光正入射在600条/mm的光栅上，第1级可见光谱末端与第2级光谱前端之间的角间隔有多少？}
\answerarea
\clearpage
\problem{4-32}{国产31W1型1米平面光栅摄谱仪的技术数据表如下，试根据这些数据来计算：
\begin{center}\begin{tabular}{l|l}\hline
物镜焦距&$1050\,\mathrm{mm}$\\光栅刻划面积&$60\,\mathrm{mm}\times40\,\mathrm{mm}$\\
闪耀波长&$635.0\,\mathrm{nm}$（1级）\\刻线&1200条/mm\\色散&$0.8\,\mathrm{nm/mm}$\\理论分辨率&72000（1级）\\\hline\end{tabular}\end{center}
（1）该摄谱仪能分辨的谱线间隔的最小值为多少？（2）该摄谱仪的角色散本领为多少〔以 $\mathrm{nm}/('')$ 为单位〕？（3）光栅的闪耀角为多大？闪耀方向与光栅平面的法线方向成多大角度？}
\answerarea
\problem{4-33}{底边长度为 $6\,\mathrm{cm}$ 的棱镜，在光波长为 $0.6\,\mu\mathrm m$ 附近能分辨的最小波长间隔为多少？设制棱镜材料的色散 $\mathrm dn/\mathrm d\lambda$ 值为 $0.4\times10^{-4}\,\mathrm{nm}^{-1}$，来估算。}
\answerarea
\clearpage
\problem{4-34}{根据以下数据比较光栅、棱镜、法--珀干涉仪三者的分光本领：（1）分辨本领；（2）色散本领；（3）自由光谱范围。
\par 光栅宽度 $D=5\,\mathrm{cm}$，刻线密度 $1/d=600$条/mm；棱镜底边 $b=5\,\mathrm{cm}$，顶角 $\alpha=60^\circ$，折射率 $n=1.51$，色散率 $\mathrm dn/\mathrm d\lambda=0.6\times10^{-4}/\mathrm{nm}$；法--珀腔长 $h=5\,\mathrm{cm}$，反射率 $R=0.99$。}
\answerarea
\clearpage

\chapter{变换光学与全息照相}
\problem{5-1}{设薄透镜由折射率为 $n_0$ 的材料做成，物方和像方的折射率分别是 $n$ 和 $n'$，导出其相位变换函数（用透镜的焦距表示出来）。}
\answerarea
\problem{5-2}{用薄透镜的相位变换函数导出傍轴条件下的横向放大率公式。}
\answerarea
\clearpage
\problem{5-3}{用劈形棱镜的相位变换函数（5.27）式导出傍轴光束斜入射时产生的偏向角 $\delta$。
\begin{equation*}
\widetilde T_p(x,y)=\exp\!\left[-\mathrm j k(n-1)\alpha x\right].\tag{5.27}
\end{equation*}
其中 $\mathrm j^2=-1$，$k$ 为入射光波数，$n$ 为棱镜折射率，$\alpha$ 为小劈角。}
\answerarea
\problem{5-4}{在一相干成像系统中，镜头（作为入射光瞳）的相对孔径为 $1/5$，求此系统的截止频率（$\mathrm{mm}^{-1}$）。设物平面在前焦面附近，照明波长为 $0.5\,\mu\mathrm m$。}
\answerarea
\clearpage
\problem{5-5}{利用阿贝成像原理导出在相干照明条件下显微镜的最小分辨距离公式，并同非相干照明的最小分辨距离公式比较。}
\answerarea
\problem{5-6}{采用远场装置（图5-20b）接收单缝的夫琅禾费衍射场，设单缝宽约为 $100\,\mu\mathrm m$，入射光波长 $632.8\,\mathrm{nm}$。
\begin{enumerate}\item 接收屏幕至少应放在多远？\item 在接收屏幕的多大范围内才算是夫琅禾费衍射场？\item 0级半角宽度为多少？\item 在接收屏幕上0级的线宽度有多少？\end{enumerate}
\bookfig[0.92]{book_5_20b}{图5-20(b)\quad 远场接收}}
\answerarea
\clearpage
\problem{5-7}{采用像面接收装置（图5-20d或e）接收单缝的夫琅禾费衍射场，设单缝宽度约为 $1\,\mathrm{mm}$，入射光波长 $488.0\,\mathrm{nm}$，物距 $40\,\mathrm{cm}$，像距 $80\,\mathrm{cm}$。
\begin{enumerate}\item 如果单缝置于透镜后方，要求在像面 $1\,\mathrm{cm}$ 范围内准确地接收到夫琅禾费衍射场，单缝离像面至少多远？\item 如果单缝紧贴透镜后面，求0级半角宽度和接收屏幕上0级的线宽度；\item 如果单缝离透镜 $40\,\mathrm{cm}$ 远，求0级半角宽度及它在幕上的线宽度；\item 如果单缝置于透镜前方，紧贴在其左侧，情况如何？\end{enumerate}
\bookfig[0.88]{book_5_20d}{图5-20(d)\quad 像面接收（一）}
\bookfig[0.88]{book_5_20e}{图5-20(e)\quad 像面接收（二）}}
\answerarea
\clearpage
\problem{5-8}{对图5-20e所示装置，推导出傍轴条件下它的衍射场表达式，并论证它符合夫琅禾费衍射场的标准形式。【提示：利用场点 $P'$ 的共轭点 $P'_c$（见图5-20e）。】
\bookfig[0.88]{book_5_20e}{图5-20(e)\quad 像面接收（二）}}
\answerarea
\problemfig[0.17]{5-9}{在透镜的前焦面上有一系列同相位的相干光源等距排列在 $x$ 轴上，形成一维点阵（见本题图），用傅里叶变换法求后焦面上的夫琅禾费衍射场。}{5_9}
\answerarea
\clearpage
\problem{5-10}{设透镜直径 $D=5\,\mathrm{cm}$，焦距 $F=60\,\mathrm{cm}$，图像（衍射屏）线度 $l=2\,\mathrm{cm}$，入射光波长 $\lambda=0.6\,\mu\mathrm m$。
\begin{enumerate}
\item 分别算出与后焦面上下列各点 $(x',y')$ 相对应的空间频率 $(f_x,f_y)$ 的具体数值：
\[
\begin{gathered}
(0,0),\quad(0,1),\quad(1,0),\quad
\left(\frac{\sqrt2}{2},\frac{\sqrt2}{2}\right),\\[4pt]
\underbrace{\left(-\frac{\sqrt2}{2}\right)}_{\text{原文坐标缺项}},\quad
(0.5,2),\quad(3,-5),\quad(-10,-15).
\end{gathered}
\]
坐标单位皆为 $\mathrm{mm}$，空间频率单位皆为 $\mathrm{mm}^{-1}$。
\item 计算系统的截止频率。
\end{enumerate}
\editorial{按补充的清晰截图核对：第五项原文仅印为 $(-\sqrt2/2)$，缺少一个坐标分量，照录并标注；其后一项明确为 $(0.5,2)$。}}
\answerarea
\clearpage
\problem{5-11}{如本题图a，参考光束 $R$ 和物光束 $O$ 均为平行光，对称地斜射于记录介质 $\Sigma$ 上，即 $\theta_R=-\theta_O$，二者间的夹角 $\theta=2\theta_O$。
\begin{enumerate}\item 试说明全息图上干涉条纹的形状；\item 试分别写出物光波和参考光波在 $\Sigma$ 平面上的相位分布 $\varphi_O(y)$ 和 $\varphi_R(y)$；\item 试证明全息图上干涉条纹的间距公式为 $d=\lambda/[2\sin(\theta/2)]$；\item 试计算，当夹角 $\theta=1^\circ$ 时间距 $d$ 为多少？当夹角 $\theta=60^\circ$ 时间距 $d$ 为多少？（采用 He--Ne 激光记录，$\lambda=632.8\,\mathrm{nm}$。）\item 某感光胶片厂生产一种可用于全息照相的记录干版，其性能如下：感光层厚度为 $8\,\mu\mathrm m$，分辨率为3000条/mm以上。利用题（4）所得数据，试说明：当夹角 $\theta=60^\circ$ 时，用该记录干版是否构成一张体全息图？当夹角 $\theta=60^\circ$ 时，该记录干版的分辨率是否匹配？\item 如本题图b，采用与参考光束 $R$ 同样波长同样倾角的照明光束 $R'$ 照射这张全息图，试分析0级、$+1$级、$-1$级三个衍射波都出现在什么方向上，并在图上画出。\end{enumerate}\fig[0.85]{5_11}
\editorial{原材料将此配图标为“习题5-12”；5-11与5-12共用此图。}}
\answerarea
\clearpage
\problem{5-12}{若在上题中改为用正入射的平面波再现，$\pm1$ 级衍射波各发生什么变化？}
\answerarea
\problem{5-13}{如本题图a，用正入射的平面参考光波记录轴外物点 $O$ 发出的球面波。
\begin{enumerate}\item 如本题图b所示，用轴上的点源 $R'$ 发出的球面波来重建波前，求 $\pm1$ 级两像点的位置；\item 用与记录全息图时不同波长的正入射平面波照明，求 $\pm1$ 级两像点的位置。\end{enumerate}\fig[0.73]{5_13}}
\answerarea
\clearpage
\problem{5-14}{（1）求图5-32a所示装置制备的全息图中各级干涉条纹的半径，并证明它们与一张菲涅耳波带片相符；（2）验证用图5-32c方式再现的两个像点 $O$、$O'$ 确是菲涅耳波带片的一对焦点。
\bookfig[0.96]{book_5_32}{图5-32\quad 菲涅耳波带片与全息图}
\par{\small 图（a）用轴上点源 $O$ 的球面波和正入射平面参考波 $R$ 记录全息图；图（b）为同心圆干涉条纹示意；图（c）用正入射平面波 $R'$ 再现，形成像点 $O$、$O'$。}}
\answerarea
\clearpage

\chapter{偏振}
\problem{6-1}{自然光投射到互相重叠的两个偏振片上。如果透射光的强度为（1）透射光束最大强度的 $1/3$；（2）入射光束强度的 $1/3$，则这两个偏振片的透振方向之间夹角是多少？假定偏振片是理想的，它把自然光的强度严格减小一半。}
\answerarea
\problem{6-2}{一束自然光入射到一偏振片组上，这组由四块组成，每片的透振方向相对于前面一片沿顺时针方向转过 $30^\circ$ 角。试问入射光中有多大一部分透过了这组偏振片？}
\answerarea
\clearpage
\problem{6-3}{将一偏振片沿 $45^\circ$ 角插入一对正交偏振器之间，自然光经过它们强度减为原来的百分之几？}
\answerarea
\problem{6-4}{计算从空气到水面的布儒斯特角（水的折射率 $n=4/3$）。}
\answerarea
\clearpage
\problem{6-5}{一束光由水射在玻璃上，当入射角为 $50.82^\circ$ 时，反射光完全偏振。求玻璃的折射率（已知水的折射率为 $4/3$）。}
\answerarea
\problem{6-6}{计算：（1）由空气到玻璃（$n=1.560$）的全偏振角；（2）由此玻璃到空气的全偏振角；（3）在全偏振时由空气到此玻璃的折射光的偏振度；（4）在全偏振时由此玻璃到空气的折射光的偏振度；（5）自然光从空气以布儒斯特角入射到平行平面玻璃板以后，最终透射光的偏振度。}
\answerarea
\clearpage
\problem{6-7}{求自然光透过八块 $n=1.560$ 的平行玻璃板组成的玻片堆后的偏振度（忽略玻璃对光的吸收）。}
\answerarea
\problemfig[0.38]{6-8}{已知自然光射于某平行平面玻璃板上时，反射光的能流为入射光的 $0.10$ 倍（见本题图），取入射能流为一个单位，设玻璃的折射率为 $1.50$，求图中标出的光束2、3、4的能流（略去玻璃对光的吸收）。}{6_8}
\answerarea
\clearpage
\problem{6-9}{线偏振光的振动面和入射面之间的夹角称为振动的方位角。设入射线偏振光的方位角为 $\alpha$，入射角为 $i_1$，求折射光和反射光的方位角 $\alpha_2$ 和 $\alpha'_1$（已知两介质的折射率为 $n_1$ 和 $n_2$）。}
\answerarea
\problem{6-10}{线偏振光以布儒斯特角从空气入射到玻璃（$n=1.560$）的表面上，其振动方位角为 $20^\circ$，求反射光和折射光的方位角。}
\answerarea
\clearpage
\problem{6-11}{设入射光、反射光、折射光的总能流分别为 $W_1$、$W'_1$、$W_2$，则总能流反射率 $\mathcal R$ 和总能流透射率 $\mathcal T$ 定义为
\[\mathcal R=\frac{W'_1}{W_1},\qquad\mathcal T=\frac{W_2}{W_1}.\]
\begin{enumerate}\item 当入射光为线偏振光，方位角（见习题6-9）为 $\alpha$ 时，试证明
\[\mathcal R=\mathcal R_p\cos^2\alpha+\mathcal R_s\sin^2\alpha,\quad\mathcal T=\mathcal T_p\cos^2\alpha+\mathcal T_s\sin^2\alpha.\]
\item 证明 $\mathcal R+\mathcal T=1$。\item 设入射光是自然光，求 $\mathcal R$、$\mathcal T$ 与 $\mathcal R_p$、$\mathcal R_s$ 和 $\mathcal T_p$、$\mathcal T_s$ 的关系。\item 设入射光是圆偏振光，求 $\mathcal R$、$\mathcal T$ 与 $\mathcal R_p$、$\mathcal R_s$ 和 $\mathcal T_p$、$\mathcal T_s$ 的关系。\end{enumerate}}
\answerarea
\clearpage
\problem{6-12}{光从空气到玻璃（$n=1.50$）以布儒斯特角入射，试计算：（1）能流反射率 $\mathcal R_p$ 和 $\mathcal R_s$ 的值；（2）能流透射率 $\mathcal T_p$ 和 $\mathcal T_s$ 的值。}
\answerarea
\problem{6-13}{线偏振光从空气到玻璃（$n=1.5$）以 $45^\circ$ 角入射，方位角为 $60^\circ$，试计算：（1）总能流反射率 $\mathcal R$ 和总能流透射率 $\mathcal T$；（2）改为自然光入射，$\mathcal R$ 和 $\mathcal T$ 为多少？}
\answerarea
\clearpage
\problemfig[0.34]{6-14}{本题图所示为一支半导体砷化镓发光管。管芯 $AB$ 为发光区，其直径 $d\sim3\,\mathrm{mm}$。为了避免全反射，发光管上部研磨成半球形，以使内部发射光能够以最大的透射率向外输出。如果要使发光区边缘两点 $A$ 和 $B$ 发射的光不至于全反射，半球的半径至少应取多少？已知砷化镓的折射率为 $3.4$（对发射的 $0.9\,\mu\mathrm m$ 波长）。}{6_14}
\answerarea
\problemfig[0.34]{6-15}{接上题，为了减少光在砷化镓--空气界面的反射，工艺上常在砷化镓表面镀一层氧化硅增透膜。如本题图所示，氧化硅的折射率为 $1.7$。现在单纯从几何光学角度提出一个问题：加膜后入射角为多少才不致于在空气表面发生全反射？试与不加膜时相比（设膜很薄，可按平面板计算）。}{6_15}
\answerarea
\clearpage
\problem{6-16}{从光密介质到光疏介质，当 $\sin i_1>n_2/n_1$ 时就发生全反射。作为一种处理方法，我们仍可在形式上维持折射定律 $n_1\sin i_1=n_2\sin i_2$，这时 $\sin i_2>1$，可认为 $i_2$ 是个虚折射角，它的余弦也为虚数：
\[\cos i_2=\sqrt{1-\sin^2 i_2}=\mathrm i\sqrt{\left(\frac{n_1}{n_2}\right)^2\sin^2 i_1-1}.\]
试利用菲涅耳公式证明此时 $|r_p|=1$、$|r_s|=1$，从而 $\mathcal R_p=\mathcal R_s=1$。
\par【注：本来这里根号前有正、负两种可能。为了保证（6.55）式中的 $\kappa>0$，从而使（6.56）式中衰减因子的指数项是负的，根号应取正的。】}
\answerarea
\problem{6-17}{推导全反射时的相移公式（6.35）。
\begin{equation*}
\left\{\begin{aligned}
\delta_p&=2\arctan\!\left[
\frac{n_1}{n_2}\frac{\sqrt{\left(\dfrac{n_1}{n_2}\right)^2\sin^2 i_1-1}}{\cos i_1}
\right],\\[7pt]
\delta_s&=2\arctan\!\left[
\frac{n_2}{n_1}\frac{\sqrt{\left(\dfrac{n_1}{n_2}\right)^2\sin^2 i_1-1}}{\cos i_1}
\right].
\end{aligned}\right.
\tag{6.35}
\end{equation*}
适用条件为 $n_1>n_2$，且 $i_1>i_c=\arcsin(n_2/n_1)$。
\editorial{原题引用（6.34），该式实际为布儒斯特角公式；依据补充的原书页面，本题应引用全反射相移公式（6.35），已据此更正。}}
\answerarea
\clearpage
\problem{6-18}{（1）计算 $n_1=1.51$、$n_2=1.0$、入射角为 $54^\circ37'$ 时全反射光的相移 $\delta_p$ 和 $\delta_s$；（2）如果入射光是线偏振的，全反射光中 p 振动和 s 振动的相位差为多少？说明由两者合成的是椭圆偏振光。}
\answerarea
\problem{6-19}{在上题中用的光源是氦氖激光，求隐失波的穿透深度。}
\answerarea
\clearpage
\problem{6-20}{一束线偏振的钠黄光垂直射入一块方解石晶体，振动方向与晶体的主平面成 $20^\circ$ 角。试计算 o、e 两光束折射光的相对振幅。}
\answerarea
\problemfig[0.39]{6-21}{两大小相同的冰洲石晶体 $A$、$B$ 前后排列，强度为 $I_0$ 的自然光垂直于晶体 $A$ 的表面入射后相继通过 $A$、$B$（见本题图）。$A$、$B$ 的主截面之间的夹角为 $\alpha$（图中 $\alpha\ne0$）。求 $\alpha=0^\circ$、$45^\circ$、$90^\circ$、$180^\circ$ 时由 $B$ 射出光线的数目和每个的强度（忽略反射、吸收等损失）。}{6_21}
\answerarea
\clearpage
\problem{6-22}{一水晶平板厚 $0.850\,\mathrm{mm}$，光轴与表面平行，用水银灯的绿光（$546.1\,\mathrm{nm}$）垂直照射，求：（1）o、e 两光束在晶体中的光程；（2）二者的相位差（用度表示）。\missing{水晶对该波长的主折射率数据表未提供。}}
\answerarea
\problem{6-23}{一束钠黄光以 $50^\circ$ 的入射角射到冰洲石平板上。设光轴与板表面平行，并垂直于入射面，求晶体中 o 光和 e 光的夹角。}
\answerarea
\clearpage
\problem{6-24}{一束钠黄光掠入射到冰的晶体平板上，光轴与入射面垂直，平板厚度为 $4.20\,\mathrm{mm}$。求 o 光和 e 光射到平板对面上两点的间隔。已知对于钠黄光冰的 $n_o=1.3090$，$n_e=1.3104$。}
\answerarea
\problem{6-25}{用 ADP（磷酸二氢铵，$\mathrm{NH_4H_2PO_4}$）晶体制成 $50^\circ$ 顶角的棱镜，光轴与折射棱平行，$n_o=1.5246$，$n_e=1.4792$。试求：（1）o 光和 e 光的最小偏向角；（2）二者之差。}
\answerarea
\clearpage
\problem{6-26}{设一水晶棱镜的顶角为 $60^\circ$，光轴平行于折射棱。钠黄光以最小偏向角的方向在棱镜中折射，用焦距为 $1\,\mathrm m$ 的透镜聚焦，o 光和 e 光两谱线的间隔为多少？\missing{水晶对钠黄光的主折射率数据表未提供。}}
\answerarea
\problem{6-27}{求冰洲石晶体中光射线和波法线间的最大夹角。}
\answerarea
\clearpage
\problem{6-28}{图6-31b中，设入射光是钠黄光，晶体为方解石，光轴与晶体表面成 $30^\circ$ 角，入射角为 $45^\circ$。（1）o 光和 e 光的方向；（2）e 光的折射率。
\bookfig[0.7]{book_6_31b}{图6-31(b)\quad 用惠更斯作图法求单轴晶体中的折射线}
\par{\small 圆形与椭圆形次波面分别对应 o 光和 e 光，图中几何形状为示意。}
\missing{方解石对钠黄光的主折射率数据；图6-31b已补入。}}
\answerarea
\problem{6-29}{当图6-46中沃拉斯顿棱镜的顶角 $\alpha=15^\circ$ 时，两出射光线间的夹角为多少？
\bookfig[0.38]{book_6_46}{图6-46\quad 沃拉斯顿棱镜}
\par{\small 两块冰洲石棱镜的光轴方向互相垂直，图中以 $\updownarrow$、$\odot$ 示意；出射方向分别标为1和2。}}
\answerarea
\clearpage
\problem{6-30}{设图6-47b所示的尼科耳棱镜中 $\angle CA''C''$ 为直角，光线 $SM$ 平行于 $A''A$，计算 $\angle S_0MS$。
\bookfig[0.94]{book_6_47b}{图6-47(b)\quad 尼科耳棱镜的主截面}}
\answerarea
\problem{6-31}{用方解石和石英薄板作对钠黄光的 $\lambda/4$ 片，它们的最小厚度各为多少？}
\answerarea
\clearpage
\problem{6-32}{两尼科耳棱镜主截面的夹角由 $30^\circ$ 变到 $45^\circ$，透射光的强度如何变化？设入射自然光的强度为 $I_0$。}
\answerarea
\problem{6-33}{单色线偏振光垂直射入方解石晶体，其振动方向与主截面成 $30^\circ$ 角，两折射光线经过置于方解石后的尼科耳棱镜，其主截面与原入射光的振动方向成 $50^\circ$ 角。求两条光线的相对强度。}
\answerarea
\clearpage
\problem{6-34}{经尼科耳棱镜观察部分偏振光，当尼科耳棱镜由对应于极大强度的位置转过 $60^\circ$ 时，光强减为一半。求光束的偏振度。}
\answerarea
\problem{6-35}{用一块 $\lambda/4$ 片和一块偏振片鉴定一束椭圆偏振光。达到消光位置时，$\lambda/4$ 片的光轴与偏振片透振方向相差 $22^\circ$，求椭圆长短轴之比。}
\answerarea
\clearpage
\problem{6-36}{强度为 $I_0$ 的右旋圆偏振光垂直通过 $\lambda/4$ 片（此 $\lambda/4$ 片由方解石做成，o 光和 e 光在晶片中的光程差刚好是 $\lambda/4$），然后再经过一块主截面相对于 $\lambda/4$ 片光轴向右旋 $15^\circ$ 的尼科耳棱镜。求最后出射的光强（忽略反射、吸收等损失）。}
\answerarea
\problem{6-37}{平行于光轴切割一块方解石晶片，放置在主截面成 $35^\circ$ 角的一对尼科耳棱镜之间，晶片的光轴平分此角。求：（1）从方解石晶片射出的 o 光和 e 光的振幅和光强；（2）由第二个尼科耳棱镜射出的 o 光和 e 光的振幅和光强。设入射自然光的光强为 $I_0=A^2$，反射和吸收等损失可忽略。}
\answerarea
\clearpage
\problem{6-38}{强度为 $I_0$ 的单色平行光通过正交尼科耳棱镜。现在两尼科耳棱镜之间插入一 $\lambda/4$ 片，其主截面与第一尼科耳棱镜的主截面成 $60^\circ$ 角。求出射光的强度（忽略反射、吸收等损失）。}
\answerarea
\problem{6-39}{一块 $0.025\,\mathrm{mm}$ 厚的方解石晶片，表面平行于光轴，放在正交尼科耳棱镜之间，晶片的主截面与它们成 $45^\circ$ 角。试问：（1）在可见光范围内哪些波长的光不能通过？（2）如果将第二个尼科耳棱镜的主截面转到与第一个平行，哪些波长的光不能通过？}
\answerarea
\clearpage
\problem{6-40}{劈形水晶棱镜顶角为 $0.5^\circ$，棱边与光轴平行，置于正交尼科耳棱镜之间，使其主截面与两尼科耳棱镜的主截面都成 $45^\circ$ 角。以水银的 $404.7\,\mathrm{nm}$ 紫色平行光正入射。
\begin{enumerate}\item 通过第二尼科耳棱镜看到的干涉图样如何？\item 相邻暗纹的间隔 $d$ 等于多少？\item 若将第二尼科耳棱镜的主截面转 $90^\circ$，干涉图样有何变化？\item 维持两尼科耳棱镜正交，但让晶片的主截面转 $45^\circ$，使之与第二尼科耳棱镜的主截面垂直，干涉图样有何变化？\end{enumerate}}
\answerarea
\clearpage
\problem{6-41}{将比比涅补偿器（见思考题6-38）放在正交偏振片之间，光轴与它们的透振方向成 $45^\circ$ 角，你将看到什么现象？若劈角 $\alpha=2.75^\circ$，用平行的钠黄光照明，求干涉条纹的间隔。转动补偿器的光轴，对干涉条纹有什么影响？
\bookfig[0.65]{review_6_38}{思考题6-38图\quad 比比涅补偿器}
\par{\small 两块劈形棱镜的光轴相互垂直，分别以点状和横线纹理示意；$\alpha$ 为劈角。}}
\answerarea
\clearpage
\problem{6-42}{以线偏振光照在比比涅补偿器上（见思考题6-38），通过偏振片观察时在中央两劈形棱镜厚度 $d_1=d_2$ 处有一暗线，与中央暗线距离 $a$ 处又有一暗线。今以一同样波长的椭圆偏振光照在此比比涅补偿器上，发现暗线移至离中央 $b$ 处。
\begin{enumerate}\item 求椭圆偏振光在补偿器晶体中分解成的两个振动分量的初始相位差与 $a$、$b$ 的关系；\item 如果椭圆的长短轴正好分别与两劈形棱镜的光轴平行，试证此时 $b=a/4$；\item 设已知偏振片的透振方向与补偿器中一劈的光轴夹角为 $\theta$，找出 $\theta$ 与（2）问中椭圆长短轴比值的关系。\end{enumerate}
\bookfig[0.6]{review_6_38}{思考题6-38图\quad 比比涅补偿器（同6-41）}}
\answerarea
\clearpage
\problem{6-43}{已知水晶对钠黄光的旋光率 $\alpha=21.75^\circ/\mathrm{mm}$，求左、右旋圆偏振光折射率之差 $\Delta n$。}
\answerarea
\problem{6-44}{在两尼科耳棱镜之间插一块石英旋光晶片，以消除对眼睛最敏感的黄绿色光（$\lambda=550.0\,\mathrm{nm}$）。设对此波长的旋光率为 $24^\circ/\mathrm{mm}$，求下列情形下晶片的厚度：（1）两尼科耳棱镜主截面正交；（2）两尼科耳棱镜主截面平行。}
\answerarea
\clearpage
\problem{6-45}{一石英棒长 $5.639\,\mathrm{cm}$，端面垂直于光轴，置于正交偏振器间，沿光轴方向输入白光，用光谱仪观察透射光。
\begin{enumerate}\item 用一大张坐标纸，画出可见光范围（$400.0$--$760.0\,\mathrm{nm}$）振动面的旋转角与波长的曲线，旋光率数据可参照表6-5；\item 从这曲线看，哪些波长的光在光谱仪中消失？\item 在这些丢失的波长中，振动面的最大和最小旋转角各是多少？\end{enumerate}\missing{原书表6-5。}}
\answerarea
\clearpage
\problem{6-46}{一块表面垂直光轴的水晶片恰好抵消 $10\,\mathrm{cm}$ 长浓度20\%的麦芽糖溶液对钠光振动面所引起的旋转。对此波长水晶的旋光率 $\alpha=21.75^\circ/\mathrm{mm}$，麦芽糖的比旋光率 $[\alpha]=144^\circ/[\mathrm{dm}\cdot(\mathrm{g/cm^3})]$。求此水晶片的厚度。}
\answerarea
\problem{6-47}{ $15\,\mathrm{cm}$ 长的左旋葡萄糖溶液使钠黄光的振动面转了 $25.6^\circ$，已知 $[\alpha]=-51.4^\circ/[\mathrm{dm}\cdot(\mathrm{g/cm^3})]$，求溶液浓度。}
\answerarea
\clearpage
\problem{6-48}{将 $14.50\,\mathrm g$ 的蔗糖溶于水，得到 $60\,\mathrm{cm^3}$ 的溶液。在 $15\,\mathrm{cm}$ 的量糖计中测得钠光振动面旋转角为向右 $16.8^\circ$。已知 $[\alpha]=66.5^\circ/[\mathrm{dm}\cdot(\mathrm{g/cm^3})]$，这蔗糖样品中有多少比例的非旋光性杂质？}
\answerarea
\problem{6-49}{钠光以最小偏向角条件射入顶角为 $60^\circ$ 的石英晶体棱镜中，棱镜中光轴与底平行。求射出的左、右旋偏振光之间的夹角（所需数据在本章给出的表格中去找）。\missing{原书本章折射率、旋光率数据表。}}
\answerarea
\clearpage

\chapter{光与物质的相互作用：光的量子性}
\problem{7-1}{有一介质，吸收系数 $\alpha=0.32\,\mathrm{cm}^{-1}$，透射光强分别为入射光强的10\%、20\%、50\%和80\%时，介质的厚度各若干？}
\answerarea
\problem{7-2}{一玻璃管长 $3.50\,\mathrm m$，内贮标准大气压下的某种气体，若这气体在此条件下的吸收系数为 $0.1650\,\mathrm m^{-1}$，求透射光强的百分比。}
\answerarea
\clearpage
\problem{7-3}{一块光学玻璃对水银灯蓝、绿谱线 $\lambda=435.8\,\mathrm{nm}$ 和 $546.1\,\mathrm{nm}$ 的折射率分别为 $1.65250$ 和 $1.62450$。用此数据定出科西公式（7.11）中的 $A$、$B$ 两常量，并用它计算对钠线 $\lambda=589.3\,\mathrm{nm}$ 的折射率 $n$ 及色散率 $\mathrm dn/\mathrm d\lambda$。}
\answerarea
\problem{7-4}{利用第一章表1-3中冕牌玻璃 K9 对 F、D、C 三条谱线的折射率数据定出科西公式（7.10）中的 $A$、$B$、$C$ 三常量，用它计算该表中给出的其它波长下折射率数据，并与表中实测数值比较。\missing{原书表1-3。}}
\answerarea
\clearpage
\problem{7-5}{一棱镜顶角 $50^\circ$，设它的玻璃材料可用二常量科西公式（7.11）来描写，其中 $A=1.53974$，$B=4.6528\times10^3\,\mathrm{nm^2}$。求此棱镜对波长 $550.0\,\mathrm{nm}$ 周围最小偏向角的色散本领。}
\answerarea
\problem{7-6}{根据（7.25）式证明吸收峰的高度反比于 $\gamma_j$，半值宽度（即峰值之半处的宽度，见图7-9b）$\Delta\lambda$ 正比于 $\gamma_j$。
\begin{equation*}
n\kappa(\lambda)=\frac{Ne^2}{2\varepsilon_0m}\frac{1}{2\pi c}
\sum_j\frac{f_j\gamma_j\lambda_j^4\lambda^3}
{(2\pi c)^2(\lambda^2-\lambda_j^2)^2+\gamma_j^2\lambda_j^4\lambda^2}.
\tag{7.25}
\end{equation*}
其中 $\lambda_j$ 为第 $j$ 类振子的共振波长，$\gamma_j$ 为阻尼参数，$f_j$ 为相应的振子强度。
\bookfig[0.88]{book_7_9}{图7-9\quad 振子的色散与共振吸收曲线}}
\answerarea
\clearpage
\problem{7-7}{一块玻璃对波长 $0.070\,\mathrm{nm}$ 的 X 射线的折射率比1小 $1.600\times10^{-6}$。求 X 射线能在此玻璃外表面发生全反射的最大掠射角。}
\answerarea
\problem{7-8}{估计一下铜的等离子体振荡角频率 $\omega_p$ 的数量级。}
\answerarea
\clearpage
\problem{7-9}{求习题7-4中冕牌玻璃 K9 对 D 双线的群速。\missing{原书表1-3。}}
\answerarea
\problem{7-10}{试计算下列各情况下的群速：
\begin{enumerate}\item $v_p=v_0$（常量）（无色散介质，如空气中的声波）。\item $v_p=\sqrt{\dfrac{\lambda}{2\pi}\left(g+\dfrac{4\pi^2T}{\lambda^2\rho}\right)}$（水面波，$g$ 为重力加速度，$T$ 为表面张力，$\rho$ 为液体的密度）。\item $n$ 满足正常色散的科西公式（7.11）。\item $\omega^2=\omega_c^2+c^2k^2$（波导中的电磁波，$\omega_c$ 为截止角频率）。\end{enumerate}}
\answerarea
\clearpage
\problem{7-11}{摄影者知道用橙黄色滤色镜拍摄天空时，在黑白照片中可增加蓝天和白云的对比。设照相机的镜头和底片的灵敏度将光谱范围限制在 $390.0\,\mathrm{nm}$ 到 $620.0\,\mathrm{nm}$ 之间，并设天空散射光在此范围内可看成是常量。若滤色镜把波长在 $550.0\,\mathrm{nm}$ 以下的光全部吸收，天空的散射光被它去掉了百分之几？}
\answerarea
\problem{7-12}{苯（$\mathrm{C_6H_6}$）的拉曼散射中较强的谱线与入射光的波数差为607、992、1178、1586、3047、$3062\,\mathrm{cm}^{-1}$。今以氩离子激光（$\lambda=488.0\,\mathrm{nm}$）入射，计算各斯托克斯和反斯托克斯光谱的波长。}
\answerarea
\clearpage
\problem{7-13}{设一个两能级系统能级差 $E_2-E_1=0.01\,\mathrm{eV}$。
\begin{enumerate}\item 分别求 $T=10^2\,\mathrm K$、$10^4\,\mathrm K$、$10^5\,\mathrm K$、$10^8\,\mathrm K$ 时粒子数 $N_2$ 与 $N_1$ 之比；\item $N_2=N_1$ 的状态相当于多高的温度？\item 粒子数发生反转的状态相当于怎样的温度？\item 我们姑且引入“负温度”的概念来描述粒子数反转的状态，你觉得 $T=-10^8\,\mathrm K$ 和 $T=+10^8\,\mathrm K$ 两个温度中哪一个更高？\end{enumerate}}
\answerarea
\clearpage

\end{document}
